\chapter{Kết Luận Và Hướng Phát Triển}

\section{Kết quả đạt được}

Đồ án đã hoàn thành một ứng dụng thuật toán đồ thị gồm hai chế độ tích hợp. Phần
định tuyến xử lý graph ô tô Hà Nội quy mô 350.930 nút và 864.982 cạnh, hỗ trợ
tối đa 10 điểm ghé, hai hàm mục tiêu và ba kiểu kết thúc. Phần thuật toán tách
rõ Dijkstra cho đường đi giữa cặp điểm, Held--Karp cho thứ tự ghé chính xác và
hai heuristic để đối chiếu. Benchmark 24 cấu hình xác nhận cả hai mặt của cách
tách này: Held--Karp giữ thời gian dưới 7,2 ms ở giới hạn 10 điểm ghé, còn
heuristic đổi bảo đảm tối ưu lấy tốc độ với cái giá đo được là tới 19,53\% chi
phí dư thừa khi tuyến bị ràng buộc điểm cuối.

Phần trò chơi mở rộng bài toán từ tối ưu tĩnh sang ra quyết định tuần tự dưới
thông tin không đầy đủ. Tập niềm tin, chu kỳ tiết lộ và ranh giới
\texttt{Observation} làm rõ máy biết gì ở mỗi lượt. Ba chiến thuật có mức phức
tạp và hành vi khác nhau, được so sánh trên 360 ván tái lập.

Bộ 112 test đều đạt trong lần chạy ngày 24/10/2026. Smoke test và kiểm tra trình
duyệt chứng minh pipeline hoạt động trên graph thật và trạng thái được giữ khi
chuyển chế độ. Sản phẩm có replay, log JSON và thông báo giới hạn dữ liệu để hỗ
trợ trình diễn và kiểm tra.

\section{Bài học thuật toán}

Thứ nhất, một bài toán ứng dụng thường cần ghép nhiều tầng tối ưu. Dijkstra đúng
cho từng đoạn chưa quyết định được thứ tự ghé; Held--Karp giải tầng tổ hợp sau khi
ma trận chi phí được tạo. Thứ hai, thuật toán chính xác và heuristic có vai trò
khác nhau: chính xác cho bảo đảm trong miền nhỏ, heuristic đổi bảo đảm lấy tốc độ
hoặc hành vi thực dụng. Thứ ba, trong hệ thống thông tin ẩn, định nghĩa giao diện
quan sát quan trọng không kém hàm đánh giá; nếu policy đọc được trạng thái thật,
mọi kết quả chiến thuật trở nên vô nghĩa.

Cuối cùng, kiểm thử cần bám vào giả định và bất biến. Kiểm tra một nút bấm hoạt
động không chứng minh thuật toán đúng; ngược lại, unit test không chứng minh giao
diện sử dụng được. Đồ án kết hợp oracle, test biên, AppTest, smoke test và trình
duyệt thật để giảm khoảng cách này.

\section{Hạn chế}

\begin{itemize}
  \item Thời gian di chuyển là tĩnh, chưa có dữ liệu ùn tắc theo thời điểm.
  \item Benchmark tầng thứ tự ghé chỉ dùng một tập 11 địa điểm và các tập con lồng
  nhau, chưa lấy mẫu ngẫu nhiên trên nhiều dạng phân bố điểm.
  \item Held--Karp tăng theo $2^n$, nên giới hạn hiện tại là 10 điểm ghé.
  \item Hệ thống chưa giải VRP nhiều xe, time window, capacity hoặc điểm động
  \cite{toth2014}.
  \item Graph trò chơi OSM không đầy đủ và được điều chỉnh hai chiều.
  \item Bắt chỉ xét vị trí nút sau bước đi; chưa xét hai tác nhân giao nhau giữa cạnh.
  \item Tập niềm tin không mang xác suất; hệ số 0,35 và shortlist 8 ứng viên chưa
  được hiệu chỉnh độc lập.
  \item Benchmark chỉ có ba bot, hai sân và 20 seed; chưa có người chơi thật hoặc
  kiểm định thống kê.
  \item Lớp tile trực tuyến phụ thuộc DNS và chính sách nhà cung cấp, dù thuật
  toán cục bộ không phụ thuộc lớp nền.
\end{itemize}

\section{Hướng phát triển có thứ tự ưu tiên}

\textbf{Ưu tiên 1 -- củng cố đánh giá.} Mở rộng seed và topology, thêm baseline
ngẫu nhiên, chạy ablation hệ số lối thoát/shortlist/chu kỳ tiết lộ, báo khoảng tin
cậy và thu thập một tập ván người chơi. Tiêu chí tiếp tục là kết quả lặp lại qua
nhiều sân; nếu ưu thế chỉ xuất hiện trên lưới, cần thu hẹp claim thay vì đổi tên
chiến thuật.

\textbf{Ưu tiên 2 -- xác suất hóa niềm tin.} Thay tập $B_t$ bằng phân phối
$p_t(v)$ với mô hình chuyển của người trốn, sau đó đánh giá expected capture hoặc
information gain. Cần có positive control: khi mô hình chuyển đúng với bot sinh
dữ liệu, phương pháp xác suất phải cải thiện so với uniform belief; nếu không,
cơ chế chưa được hỗ trợ.

\textbf{Ưu tiên 3 -- lập kế hoạch nhiều bước.} So sánh heuristic một bước với
beam search hoặc Monte Carlo tree search \cite{browne2012} trên sân nhỏ. Trước khi mở rộng, cần đo
trade-off giữa tỉ lệ bắt và độ trễ; không gọi là tối ưu nếu chưa có chứng minh
hoặc oracle kích thước nhỏ.

\textbf{Ưu tiên 4 -- định tuyến động và bền vững.} Nếu có dữ liệu thời gian theo
cạnh $c_{ij}(t)$, mô hình phải kiểm tra thuộc tính FIFO và chạy thuật toán phù
hợp; không dùng một hệ số ``giờ cao điểm'' toàn cục rồi gọi đó là time-dependent
routing. Một hướng gần hơn là cho phép chặn một số cạnh và đo độ tăng chi phí,
tạo kịch bản ``phá đường--cứu tuyến'' có đối chứng rõ ràng.

\section{Kết luận chung}

Hệ thống cho thấy một đề tài tối ưu lộ trình truyền thống có thể phát triển thành
một phòng thí nghiệm thuật toán đồ thị có tính tương tác mà vẫn giữ được phần lõi
ban đầu. Giá trị chính không nằm ở tuyên bố một chiến thuật luôn tốt nhất, mà ở
việc mô hình hóa rõ, triển khai có ranh giới thông tin, so sánh đúng mức và cung
cấp bằng chứng tái lập. Đây là nền tảng phù hợp để tiếp tục nghiên cứu các bài
toán định tuyến và pursuit--evasion phức tạp hơn.
